% TOP_PROBLEM: {"problemId": "problem.zilber-pink-conjecture-on-unlikely-intersections", "canonicalTitle": "Zilber-Pink Conjecture on Unlikely Intersections", "releaseRank": 82, "primaryDomain": "number_theory_arithmetic_geometry", "primaryDomainLabel": "Number theory & arithmetic geometry", "displayStatus": "open_with_solved_subcases", "releaseStatus": "open_with_solved_subcases"}
% !TeX program = lualatex
% ATTEMPT_STATUS: unresolved
\documentclass[11pt]{article}
\usepackage[margin=25mm]{geometry}
\usepackage{fontspec}
\setmainfont{Latin Modern Roman}
\usepackage{amsmath,amssymb,mathtools}
\usepackage{unicode-math}
\setmathfont{Latin Modern Math}
\usepackage{xurl,hyperref}
\hypersetup{colorlinks=true,urlcolor=blue}
\setcounter{secnumdepth}{0}
\setlength{\parindent}{0pt}
\setlength{\parskip}{5pt}
\setlength{\emergencystretch}{3em}
\begin{document}
\section*{\texorpdfstring{Zilber-Pink Conjecture on Unlikely Intersections}{Zilber-Pink Conjecture on Unlikely Intersections}}\label{82}
\subsection{Definitions and mathematical statement}
A special subvariety of a mixed Shimura variety is one arising from the appropriate Shimura subdatum and Hecke-translate construction. For irreducible $W\subseteq S$, the expected intersection dimension with $T$ is $\dim W+\dim T-\dim S$. The record calls an irreducible $A\subseteq W\cap T$ atypical if
\[
 \dim A>\dim W+\dim T-\dim S.
\]
The claim is that only finitely many such $A$ are maximal under inclusion. Equivalently their union is the union of finitely many atypical subvarieties. The source's special-subvariety convention is essential; replacing ``special'' with an arbitrary algebraic family would give a different problem.
\par\smallskip\textit{Formulation and concept references:} \href{https://ems.press/content/serial-article-files/52570}{[S1]}.

\subsection{Short English statement}
Are all unexpectedly large special intersections inside a fixed Shimura subvariety controlled by finitely many maximal exceptional pieces?
\subsection{Research attempt}
\paragraph{Process, scope, and first gap.}
This attempt was produced by GPT-6 Astra Ultra on 27 September 2026, by checking the cited formulation, reducing the dimension inequality, and testing an attempted finiteness argument against explicit algebraic examples. The calculations below are elementary reductions and applications of named theorems, with no novelty claim. Write $s=\dim S$ and $w=\dim W$. Subvarieties are closed and irreducible unless otherwise indicated. Work in the relevant irreducible component of $S$.

The first unproved assertion is already concrete for a curve $W$ contained in no proper special subvariety of $S$:
\[
 \#\left(W\cap\bigcup_{\substack{T\text{ special in }S\\
                         \operatorname{codim}_S T\ge2}}T\right)<\infty.
 \tag{82.1}
\]
For arbitrary mixed Shimura $S$, this attempt does not prove (82.1). In higher dimension one must also control maximal positive-dimensional atypical pieces. The objective is finiteness for each fixed $W$, with no asserted uniform bound across all $W$.

\paragraph{Literature and the arithmetic input.}
The ambient-dimension convention is that of [S1, Conjecture 2.3]. The stronger hypotheses in the non-$\overline{\mathbb Q}$ results of [S2] cannot be discarded to obtain an assertion for every algebraic $W$. One established input used below is Pila's Andr\'e--Oort theorem for products of modular curves [E1, Theorem 1.1*]: the irreducible components of the Zariski closure of a set of special points are special. The theorem applies to subvarieties defined over $\mathbb C$; its special points and special subvarieties have their modular meaning. Only the modular-surface consequence is used here. No general mixed-Shimura theorem is inferred from that case.

\paragraph{Replacing arbitrary pieces by components.}
Suppose $A\subset W\cap T$ satisfies the strict inequality in the statement. Choose an irreducible component $B$ of $W\cap T$ containing $A$. Then
\[
 \dim B\ge\dim A>w+\dim T-s,
\]
so $B$ is atypical too. Consequently every maximal atypical $A$ is an actual component of some intersection. Moreover every atypical $A$ is contained in a maximal atypical subvariety: among atypical subvarieties containing $A$, choose one of largest dimension. A proper containment between irreducible closed subvarieties strictly increases dimension, so this choice is maximal.

This gives the equivalence of the two finiteness formulations without taking the Zariski closure of the exceptional union. If finitely many atypical $A_i$ cover that union, an irreducible maximal atypical $B$ is contained in one $A_i$, because an irreducible space cannot be a finite union of proper closed subsets. Maximality then gives $B=A_i$. Conversely finitely many maximal pieces cover all atypical pieces by the preceding dimension argument. The bounded dimension proves existence of maximal elements; it does not bound how many there are.

For any fixed finite list $T_1,\ldots,T_m$, the assertion holds immediately. Each $W\cap T_i$ has finitely many irreducible components. Every atypical piece witnessed by this list lies in one of the finitely many atypical components. Thus the missing issue is uniform control as the special subvariety varies through the entire infinite collection.

\paragraph{Two boundary reductions and a codimension test.}
If $W$ is contained in a proper special $T_0$, then
\[
 w>w+\dim T_0-s.
\]
Hence $W$ itself is atypical and is the unique maximal atypical piece. This case is genuinely trivial for the recorded ambient convention. It would be a mistake to silently replace $S$ by $T_0$ and claim the resulting, usually stronger, question had been solved. If $W=S$, on the other hand, any $A\subset T$ has $\dim A\le\dim T$, and there are no atypical pieces.

For a proper $A\subset W$, set $r=w-\dim A$. The atypicality test becomes the exact integer inequality
\[
 \operatorname{codim}_S T>r.
 \tag{82.2}
\]
Thus an atypical divisor of $W$ must be witnessed by a special subvariety of ambient codimension at least two. An atypical point requires ambient codimension at least $w+1$. In particular, intersecting a curve with a special divisor at a point has the expected dimension zero and is not atypical merely because the point is arithmetically interesting. The strict inequality is essential.

When $w=0$, there is at most one maximal piece, namely $W$ itself if it is atypical. When $s=1$, irreducible $W$ has dimension zero or is the whole component, so these boundary arguments settle that ambient dimension completely.

\paragraph{Exact reduction for curves.}
Assume $w=1$ and that $W$ is contained in no proper special subvariety. A positive-dimensional irreducible closed $A\subset W$ equals $W$. It cannot be atypical: a witness $T$ would have to contain $W$ and be proper in $S$. Every atypical piece is therefore a point.

For a point $a\in W\cap T$, its atypicality is
\[
 0>1+\dim T-s,
\]
which is precisely $\operatorname{codim}_S T\ge2$. Conversely every point in (82.1) satisfies this inequality. Distinct points do not contain one another, so each is maximal. This proves that (82.1) is exactly the original finiteness assertion in this case, rather than just a necessary condition. It also identifies why a proof that controls only positive-dimensional intersections leaves the first difficult case untouched.

\paragraph{A completed modular-surface subcase.}
Take $S=Y(1)^2$, the product of two coarse modular curves, and an arbitrary irreducible algebraic $W\subset S$. If $W$ has dimension zero or two, use the boundary reductions. If $W$ is a curve contained in a proper special subvariety, $W$ itself is atypical as above.

In the remaining curve case, a special $T$ of codimension at least two is a special point: a pair of complex-multiplication $j$-invariants. Were infinitely many such points to lie on $W$, they would be Zariski dense in the irreducible curve. The theorem from [E1] would then make $W$ special, contrary to the assumption. Therefore the set is finite, and the full recorded assertion holds in this modular surface. The elementary work is the reduction to the exact theorem hypothesis; the finiteness theorem is an external input, not established anew here.

The same argument does not extend by simply writing $Y(1)^n$ in place of $Y(1)^2$. For $n\ge3$, codimension-two special varieties can have positive dimension. In $Y(1)^3$, for example,
\[
 Y(1)\times\{c_2\}\times\{c_3\}
\]
is special when $c_2,c_3$ are CM values. Its points with non-CM first coordinate are not special points of the product. A point of intersection with a curve nevertheless has the strict dimension excess required by (82.2). A special-point theorem therefore does not cover all the intersections appearing in the curve reduction.

\paragraph{A failed Noetherian argument, with a countertest.}
A tempting proposal is to enumerate special $T_i$, form the finite union of atypical components obtained up to step $m$, and appeal to Noetherianity to conclude that this increasing sequence of closed sets stabilizes. Noetherianity gives stabilization for decreasing chains of closed sets, not increasing ones.

The failure can be made algebraic and dimensionally exact. Let $S=\mathbb A^2_{\mathbb C}$, $W=\{y=0\}$, and choose the family
\[
 T_n=\{(n,0)\},\qquad n\in\mathbb Z_{\ge1}.
\]
Relative to this deliberately arbitrary family, every $T_n$ is an atypical point of $W$, since $0>1+0-2=-1$. The union of the first $m$ points is closed and strictly increases with $m$. Their infinite union is dense in $W$, because a nonzero univariate polynomial has only finitely many roots, but is not closed. None of the individual witnesses contains $W$; the closure $W$ is not an atypical piece for this family.

This is not a counterexample to Zilber--Pink: these $T_n$ are not being asserted to be Shimura-special. It isolates the exact invalid step in an argument that uses only finite-dimensional algebraic geometry. The same example shows that bounding the degree and number of components of each separate intersection would not suffice: here each intersection has degree one and a single component.

\paragraph{Verification and next obstruction.}
The dimension calculations were checked at codimensions zero, one, and two and at $w=0,1,s$. Empty intersections contribute no irreducible pieces. The component argument never takes the closure of the infinite exceptional union. The modular result uses the theorem only where every remaining witness is zero-dimensional; it does not identify all atypical points with special points in higher dimension. These are symbolic checks, not a computation over CM values or an exhaustive verification of Shimura subvarieties.

A useful next target is an arithmetic complexity bound on the codimension-at-least-two witnesses meeting a fixed curve, strong enough to imply (82.1). A second target is a finite collection of proper atypical loci containing every positive-dimensional exceptional component for a fixed higher-dimensional $W$. Neither target follows from degree bounds on one intersection or from bounded lengths of inclusion chains. A potential counterexample would need infinitely many incomparable maximal pieces satisfying the actual special-subvariety conditions, rather than the arbitrary points in the countertest.

\paragraph{Outcome.}
\textbf{UNRESOLVED.} The component and codimension reductions are proved, the modular-surface case follows from the stated Andr\'e--Oort input, and the proposed purely Noetherian route is refuted. The finiteness assertion for a general mixed Shimura variety remains unproved.

\subsection{Sources}
\begin{itemize}
\item[S1]J. Pila, 'Combining the conjectures of Schanuel and Zilber-Pink', Rend. Lincei Mat. Appl. 36 (2025), 653-670, DOI 10.4171/RLM/1086; accepted 1 July 2025.
\url{https://ems.press/content/serial-article-files/52570}.
\textit{Formulation source.}
\item[S2]The Zilber-Pink conjecture for varieties not defined over Qbar.
\url{https://arxiv.org/html/2504.00865v1}.
\textit{Further reference; not a proof endorsement.}
\item[S3]A note on unlikely intersections in Shimura varieties.
\url{https://arxiv.org/html/2209.07967}.
\textit{Further reference; not a proof endorsement.}
\item[E1]J. Pila, \emph{O-minimality and the Andr\'e--Oort conjecture for $\mathbb C^n$}, Annals of Mathematics 173 (2011), 1779--1840, Theorems 1.1 and 1.1*.
\url{https://annals.math.princeton.edu/wp-content/uploads/annals-v173-n3-p11-p.pdf}.
\textit{Established input for products of modular curves; consulted 27 September 2026.}
\end{itemize}
\end{document}
